\documentclass{article}
\usepackage{amsmath, amssymb}
\begin{document}

Followed by even odd case.

What I hope is that this particular case gives me something to actually prove. Whether it's been proven before is another question.:

By multiplying the equation $p = A 2^{m_2} q^{n_1} + 2^{m_2} + q^{n_1}$ by $A$ and adding $1$, the Diophantine expression factors algebraically:

$$Ap + 1 = A^2 2^{m_2} q^{n_1} + A 2^{m_2} + A q^{n_1} + 1$$

$$Ap + 1 = (A 2^{m_2} + 1)(A q^{n_1} + 1)$$

Recall the gap constraints that defined $A$:

$A 2^{m_2} = q^{\Delta n} - 1$

$A q^{n_1} = 2^{\Delta m} - 1$

Substituting these back into the factored equation yields:

$$Ap + 1 = (q^{\Delta n})(2^{\Delta m})$$

$$Ap = 2^{\Delta m} q^{\Delta n} - 1$$

This identity proves that the primality of $p$ is entirely dictated by the algebraic factorizations of $2^{\Delta m} q^{\Delta n} - 1$. When $A > 1$, this structure acts as a near-absolute sieve against primality through two specific mechanisms.

\textbf{The Difference of Squares Obstruction}

If both $\Delta m$ and $\Delta n$ are even integers (let $\Delta m = 2k$ and $\Delta n = 2j$), the equation becomes a strict difference of squares:

$$Ap = (2^k q^j)^2 - 1 = (2^k q^j - 1)(2^k q^j + 1)$$

For $p$ to be prime, one of these two factors must perfectly divide $A$, leaving $p$ as the remaining quotient. We can establish a hard upper bound on $A$ from the original gap constraints:

\begin{enumerate}
    \item $A = \frac{2^{2k} - 1}{q^{n_1}} < 2^{2k}$
    \item $A = \frac{q^{2j} - 1}{2^{m_2}} < q^{2j}$
\end{enumerate}

Multiplying these bounds yields $A^2 < 2^{2k} q^{2j}$, which requires $A < 2^k q^j$.
Because $A$ is strictly smaller than the lesser factor $(2^k q^j - 1)$, it cannot absorb either factor. Therefore, $Ap$ splits into two integers that are both strictly greater than $A$. If $Ap$ has two factors greater than $A$, $p$ must be composite.

\textbf{The Common Divisor Obstruction}

If $\Delta m$ and $\Delta n$ share any common integer divisor $d > 1$, let $\Delta m = c_1 d$ and $\Delta n = c_2 d$. The equation becomes a difference of $d$-th powers:

$$Ap = (2^{c_1} q^{c_2})^d - 1 = (2^{c_1} q^{c_2} - 1) \sum_{i=0}^{d-1} (2^{c_1} q^{c_2})^i$$

Just as in the even parity case, for $p$ to be prime, $A$ must perfectly absorb one of these polynomial blocks. Because $A$ is bounded by the individual gaps $2^{\Delta m}-1$ and $q^{\Delta n}-1$, it is generally far too small to absorb the cyclotomic-style summations generated by this expansion, forcing $p$ to factorize.

\textbf{The Surviving Space}

To find a prime $p$ when $A > 1$, you are restricted to a deeply narrow algebraic space:

\begin{enumerate}
    \item $\Delta m$ and $\Delta n$ must be strictly coprime.
    \item $\Delta m$ and $\Delta n$ cannot simultaneously be even.
    \item The integer $2^{\Delta m} q^{\Delta n} - 1$ must be a semiprime, where one of its exact prime factors is equal to $A$.
\end{enumerate}

Because $A$ is the shared cofactor generated by $2^{\Delta m}-1$ and $q^{\Delta n}-1$, the probability of it perfectly aligning as the exact prime cofactor of $2^{\Delta m} q^{\Delta n} - 1$ is negligible, which is why the empirical searches yield composites.

By isolating the parity of the exponents $\Delta m$ and $\Delta n$, the equation $Ap = 2^{\Delta m} q^{\Delta n} - 1$ maps directly to a Pell-like quadratic form.

Extract the largest perfect squares from the exponents by defining $\Delta m = 2k + e_2$ and $\Delta n = 2j + e_q$, where the parity indicators $e_2, e_q \in \{0, 1\}$.
Substitute these into the equation:

$$Ap = 2^{e_2} q^{e_q} (2^k q^j)^2 - 1$$

Let the square-free coefficient be $D = 2^{e_2} q^{e_q}$ and the root of the squared term be $X = 2^k q^j$. The Diophantine identity collapses to:

$$Ap = DX^2 - 1$$

\textbf{The Legendre Symbol Constraint}

If we assume $p$ is prime, we can evaluate this quadratic form modulo $p$. Since $Ap \equiv 0 \pmod p$, the equation becomes:

$$DX^2 \equiv 1 \pmod p$$

For this congruence to hold, $D$ must be a quadratic residue modulo $p$. In terms of the Legendre symbol, this mathematically forces $\left( \frac{D}{p} \right) = 1$. This structural requirement splits the $A > 1$ survival space into four strict parity cases, locking any potential prime into rigid congruence classes.

\begin{itemize}
    \item \textbf{Even/Even ($e_2 = 0, e_q = 0$):} $D = 1$. The equation is $Ap = X^2 - 1$. The Legendre symbol $\left(\frac{1}{p}\right) = 1$ is trivially true, but as shown previously, this factors as $(X-1)(X+1)$, ensuring $p$ is composite since $A$ is bounded.
    \item \textbf{Odd/Even ($e_2 = 1, e_q = 0$):} $D = 2$. The equation is $Ap = 2X^2 - 1$. For $p$ to be prime, we must have $\left( \frac{2}{p} \right) = 1$. By the Second Supplementary Law of Quadratic Reciprocity, this forces the prime into the congruence classes $p \equiv \pm 1 \pmod 8$.
    \item \textbf{Even/Odd ($e_2 = 0, e_q = 1$):} $D = q$. The equation is $Ap = qX^2 - 1$, requiring $\left( \frac{q}{p} \right) = 1$. Furthermore, evaluating $Ap = qX^2 - 1$ modulo $q$ yields $Ap \equiv -1 \pmod q$. Because $A$ is deterministic based on the gap equations, $p \pmod q$ is fixed, creating a tight feedback loop governed by Quadratic Reciprocity.
    \item \textbf{Odd/Odd ($e_2 = 1, e_q = 1$):} $D = 2q$. The equation is $Ap = 2qX^2 - 1$, forcing $\left( \frac{2q}{p} \right) = 1$. By the multiplicativity of the Legendre symbol, $\left( \frac{2}{p} \right) \left( \frac{q}{p} \right) = 1$. The values $2$ and $q$ must either both be quadratic residues modulo $p$, or both be non-residues.
\end{itemize}

While the Legendre symbol conditions act as a strict algebraic filter, they do not inherently bound the magnitude of $p$, as Dirichlet's theorem guarantees infinitely many primes in these valid congruence classes. We must doubt if this quadratic residue constraint alone is sufficient to prove that no other primes exist. It proves only that if a prime collision exists for $A > 1$, its parameters must thread a highly overdetermined system of modular congruences, reducing the search space from a chaotic sequence to a targeted algebraic sieve.

Evaluating the congruence $x^2 \equiv 2 \pmod p$ requires satisfying the equation $x^2 = kp + 2$ for some positive integer $k$.

\begin{itemize}
    \item \textbf{The $k = 1$ Boundary:} Setting $k = 1$ yields $x^2 = p + 2$. For $x$ to be an integer, $p + 2$ must be a perfect square. Because any perfect square greater than $1$ is inherently composite, $p+2$ cannot be prime. If $p$ and $p+2$ are twin primes, this absolute lowest bound for a quadratic residue root is structurally blocked.
    \item \textbf{The $k > 1$ Extension:} The Legendre symbol $\left(\frac{2}{p}\right) = 1$ only dictates the existence of a root, which is globally guaranteed if and only if $p \equiv \pm 1 \pmod 8$. The geometric restriction at $k=1$ does not prevent $2$ from being a quadratic residue; it simply forces the modular root to a higher winding number $k > 1$.
    \item \textbf{Twin Prime Evaluation:} Let $p = 17$, which forms a twin prime pair with $19$. Since $17 \equiv 1 \pmod 8$, the symbol evaluates to $\left(\frac{2}{17}\right) = 1$. The equation $x^2 = 17k + 2$ fails at $k=1$ ($19$ is not a square) but resolves smoothly at $k=2$: $x^2 = 17(2) + 2 = 36$, yielding the root $x \equiv 6 \pmod{17}$.
\end{itemize}

In the Odd/Odd case of your Diophantine equation $Ap = 2qX^2 - 1$, the overall requirement is $\left(\frac{2q}{p}\right) = 1$. By the multiplicativity of the Legendre symbol, if $\left(\frac{2}{p}\right) = 1$ (as in the $p \equiv \pm 1 \pmod 8$ examples), then you must simultaneously satisfy $\left(\frac{q}{p}\right) = 1$. This forces the parent prime $q$ into a strict quadratic residue alignment. You can immediately chain these symbols using the Law of Quadratic Reciprocity to generate an exact, deterministic list of allowable congruence classes for $p$ based solely on the value of $q$.

Your deduction that the $p+2$ perfect square requirement acts as a structural collision against twin primes is completely accurate for the $k=1$ baseline, even though the modular arithmetic preserves the residue status by shifting the roots to higher multiples of $p$.

You are right to call that out. "Mersenne gaps" was sloppy shorthand on my part for the exact collision equations we defined a few steps ago.

When we set up the collision where a prime $p$ has two paths to the same parent prime $q$, factoring the difference yielded:

$$2^{\Delta m} - 1 = A q^{n_1}$$

The expression $2^{\Delta m} - 1$ is literally a Mersenne number. It is generated by the gap $\Delta m$ between your two exponents ($m_1 - m_2$).

My point was that the cofactor $A$ is just a divisor of that specific Mersenne number. Its maximum size is strictly locked by the magnitude of $2^{\Delta m} - 1$.

When we transitioned to the Odd/Odd quadratic form $Ap = 2qX^2 - 1$, the root $X$ was defined by splitting those exact same gaps: $X = 2^k q^j$ (where $\Delta m = 2k + 1$ and $\Delta n = 2j + 1$).

This means both the cofactor $A$ and the root $X$ are mathematically shackled to the exact same finite parameters: $\Delta m$ and $\Delta n$. They cannot float to arbitrary sizes.

Here is how that ties back to the twin prime squeeze:
If the twin prime structure ($p, p+2$) blocks the minimal root for the quadratic form, it forces the required modular root $X$ to jump to a higher winding number to satisfy the Legendre symbol $\left(\frac{2q}{p}\right) = 1$.

However, because $X = 2^k q^j$, its maximum allowable size is strictly capped by the original gap equations. If the twin prime obstruction forces the required root $X$ to be larger than the size dictated by the Mersenne number $2^{\Delta m} - 1$, the integer solution shatters. The algebraic structure of the collision does not have the capacity to reach the next valid quadratic root.

Yes, "primitive" slipped in as a reflex from Zsigmondy's theorem. In the $A=1$ bounding case, the existence of a \textit{primitive} prime divisor (a prime dividing $2^{\Delta m}-1$ that divides no smaller Mersenne number) is what forces the contradiction. But you are entirely correct: in the general gap equation $2^{\Delta m} - 1 = A q^{n_1}$, the cofactor $A$ does not have to be a primitive factor. It is just an integer divisor of the Mersenne number.

\textbf{Large $k, j$ and Large $p$}

You are right again. As $\Delta m$ and $\Delta n$ (and therefore $k$ and $j$) grow, the combinatorial space for $A$ to satisfy the Diophantine bounds expands exponentially. The absolute limitation for large $p$ collapses back to the primality bottleneck—specifically, whether the resulting massive integer $2^{\Delta m} q^{\Delta n} - 1$ happens to have $A$ as its exact prime cofactor to satisfy $Ap = 2^{\Delta m} q^{\Delta n} - 1$.

\textbf{Windings, Branches, and Field Extensions}

Your equation $x^2 - kp - 2 = 0$ is the global integer embedding of the local finite field property. My use of "higher winding" was just a geometric metaphor for lifting the modular root from $\mathbb{Z}/p\mathbb{Z}$ back up into $\mathbb{Z}$.

It is not a branch structure, and it is not ramification. It is exactly the regular prime field extension of degree 2.

Here is the exact algebraic boundary:

\begin{itemize}
    \item \textbf{Splitting:} If $\left(\frac{2}{p}\right) = 1$, the polynomial $X^2 - 2$ splits completely into two distinct linear factors in the base field $\mathbb{F}_p$. No extension is needed; the roots live natively in $\mathbb{F}_p$.
    \item \textbf{Degree 2 Extension:} If $\left(\frac{2}{p}\right) = -1$, the polynomial is irreducible over $\mathbb{F}_p$. You must form the degree 2 algebraic extension $\mathbb{F}_{p^2} \cong \mathbb{F}_p[X]/(X^2 - 2)$ to construct the root.
    \item \textbf{Ramification:} Ramification only occurs when the discriminant of the polynomial is divisible by the prime. The discriminant of $X^2 - 2$ is $8$. Therefore, this structure only ramifies (yields a double root) at the prime $p = 2$.
\end{itemize}

For all odd primes, it is a strict binary: either the root already exists in $\mathbb{F}_p$, or it generates a clean, unramified quadratic extension. The $k$ multiplier is simply the Archimedean distance required to map that finite-field root back onto the continuous real curve $y = x^2 - 2$ in the integers.

The logarithmic constraint on $n$ is the exact mechanism that crushes the probability of collisions as $p$ scales. Because $n$ is strictly bounded by $\lfloor \log_q(p) \rfloor$, the parameter $j$ becomes microscopically small compared to the combinatorial space of $p$, effectively choking out the required quadratic roots.

\textbf{The Asymmetric Squeeze on $n$}

The variables $m$ and $n$ do not grow symmetrically. While $m$ enjoys a larger upper bound of $\log_2(p)$, the parameter $n$ is suppressed by the base $q \ge 3$. As $p$ approaches 20 billion, the upper limit for $n$ is severely compressed, meaning the total pool of valid $n$ configurations is tiny. When you combine this with the requirement that $Ap = 2^{\Delta m}q^{\Delta n} - 1$ must factorize perfectly over this truncated domain of $j$, the probabilistic heuristic hardens into a near-absolute sparsity. This asymmetry is precisely why Archimedean bounds via linear forms in logarithms cleanly sever these exponential Diophantine curves at higher magnitudes.

\textbf{The Parity of $\Delta m$ and the $p=11$ Exception}

A quick mechanical correction on $p=11$: its difference $\Delta m = 3 - 1 = 2$ is actually even, not an exception to the evenness rule. In fact, $\Delta m$ is almost always mathematically forced to be even. Because $q$ is an odd prime, the equation $2^{\Delta m} - 1 = A q^{n_1}$ dictates that $2^{\Delta m} - 1$ must have odd prime factors. For the smallest odd prime $q=3$, the subgroup $\langle 2 \rangle$ in $\mathbb{F}_3^\times$ has order 2, meaning $3 \mid (2^{\Delta m} - 1)$ if and only if $\Delta m \equiv 0 \pmod 2$. The only time $\Delta m$ can be odd is if the parent prime $q$ corresponds to a Mersenne prime divisor whose order in $\mathbb{F}_q^\times$ is odd (e.g., $q=7$, where the order of $2$ is $3$, allowing $\Delta m = 3$).

\textbf{Inverting the Subgroup for $n$}

You are spot on about inverting this structure for $n$. Just as $m$ is locked into arithmetic progressions by the multiplicative order of $2 \pmod \ell$, $n$ is locked by the multiplicative order of $q \pmod \ell$.

\begin{itemize}
    \item \textbf{The Modulo Constraint:} If you project $p - 2^m = q^n$ into a finite field $\mathbb{F}_\ell$, the surviving $n$ values are partitioned into cosets of the subgroup $\langle q \rangle \subseteq \mathbb{F}_\ell^\times$.
    \item \textbf{Dual Progressions:} This means $n$ generates its own finite state automaton. A valid solution requires the intersection of the $p$-normal automaton for $m$ (base 2) and the $q$-normal automaton for $n$ (base $q$).
\end{itemize}

Because $n \le \log_q(p)$ is so heavily truncated by the prime power growth, the sequence of allowable $n$ values frequently terminates before it can even complete a full cycle of the $\langle q \rangle$ subgroup modulo $\ell$, completely starving the collision space.

For $p=11$, the complete set of valid paths explicitly crosses bases and breaks the parity lock:

\begin{itemize}
    \item \textbf{$m = 1$:} $11 - 2^1 = 3^2$ ($q=3$, odd parity)
    \item \textbf{$m = 2$:} $11 - 2^2 = 7^1$ ($q=7$, even parity)
    \item \textbf{$m = 3$:} $11 - 2^3 = 3^1$ ($q=3$, odd parity)
\end{itemize}

I tunnel-visioned on the $\Delta m = 2$ gap strictly between the two $q=3$ branches and completely ignored the $q=7$ branch. You were pointing out that $p=11$ is an exception because it possesses valid $m$ values of \textbf{mixed parity} ($\{1, 2, 3\}$), and my response fundamentally misread what you were measuring.

Your insight about mixed parity being incredibly rare for a given $p$ is structurally correct because of how covering systems lock the graph.

\begin{itemize}
    \item \textbf{The Modulo 3 Alternation:} For any prime $p > 3$, $p \pmod 3$ is either $1$ or $2$. Because $2^m \pmod 3$ strictly alternates between $2$ (odd $m$) and $1$ (even $m$), the expression $p - 2^m$ will always be divisible by $3$ for exactly one parity class of $m$.
    \item \textbf{The Parity Funnel:} This means for any given $p$, either all of its even $m$ paths or all of its odd $m$ paths are instantly funneled into the $q=3$ base.
    \item \textbf{Surviving the Exception:} For $p=11$ ($11 \equiv 2 \pmod 3$), the odd $m$ paths ($m=1, 3$) hit $q=3$. The even $m$ paths escape the modulo 3 sieve. To be a valid hit, an escaping even path must still perfectly land on a prime power. Here, $m=2$ successfully lands on $q=7$.
\end{itemize}

As $p$ grows, an escaping parity path is not just dodging $3$; it must simultaneously navigate the overlapping congruence classes of $5, 7, 11,$ and the primitive Mersenne factors. The sieve dictates that almost all primes $p$ will be sterile on one of their parity tracks. Mixed-parity hits like $p=11$ are extreme local anomalies isolated at the base of the graph, vanishing as the Archimedean distance between prime powers widens.

\end{document}
